\small
\begin{enumerate}[label=\textbf{\arabic*.}, leftmargin=0pt, itemsep=1.2em]
\item Какую основную архитектуру имеет двухслойная нейронная сеть (Многослойный персептрон, MLP)?
\begin{itemize}[label=$\square$, leftmargin=1.5cm, itemsep=0.2em, topsep=0.3em]
\item Только входной и выходной слои
\item Входной слой, один скрытый слой и выходной слой
\item Сеть без слоев, полносвязный граф случайной структуры
\item Только два параллельных входных слоя
\end{itemize}
{\small\textit{Правильный ответ: Входной слой, один скрытый слой и выходной слой.}}

\item Какая теорема утверждает, что нейронная сеть с одним скрытым слоем может аппроксимировать любую непрерывную функцию с заданной точностью?
\begin{itemize}[label=$\square$, leftmargin=1.5cm, itemsep=0.2em, topsep=0.3em]
\item Теорема Пифагора
\item Универсальная теорема аппроксимации (Цыбенко или Хорника)
\item Теорема Байеса
\item Теорема Черча-Тьюринга
\end{itemize}
{\small\textit{Правильный ответ: Универсальная теорема аппроксимации (Цыбенко или Хорника).}}

\item Зачем в скрытом слое нейронной сети используются нелинейные функции активации (например, Сигмоида или ReLU)?
\begin{itemize}[label=$\square$, leftmargin=1.5cm, itemsep=0.2em, topsep=0.3em]
\item Чтобы уменьшить количество весов
\item Потому что композиция линейных функций остается линейной, и без нелинейности скрытые слои бесполезны
\item Для ускорения вычислений
\item Чтобы веса не становились отрицательными
\end{itemize}
{\small\textit{Правильный ответ: Потому что композиция линейных функций остается линейной, и без нелинейности скрытые слои бесполезны.}}

\item Какой алгоритм используется для обучения многослойных нейронных сетей?
\begin{itemize}[label=$\square$, leftmargin=1.5cm, itemsep=0.2em, topsep=0.3em]
\item Генетический поиск
\item Метод k-ближайших соседей
\item Алгоритм обратного распространения ошибки (Backpropagation)
\item Случайный лес
\end{itemize}
{\small\textit{Правильный ответ: Алгоритм обратного распространения ошибки (Backpropagation).}}

\item На каком математическом правиле основан алгоритм обратного распространения ошибки?
\begin{itemize}[label=$\square$, leftmargin=1.5cm, itemsep=0.2em, topsep=0.3em]
\item Правило дифференцирования сложной функции (Chain rule)
\item Формула Тейлора
\item Преобразование Фурье
\item Ряд Маклорена
\end{itemize}
{\small\textit{Правильный ответ: Правило дифференцирования сложной функции (Chain rule).}}

\item Что такое прямое распространение (forward propagation)?
\begin{itemize}[label=$\square$, leftmargin=1.5cm, itemsep=0.2em, topsep=0.3em]
\item Процесс обновления весов от выхода ко входу
\item Процесс последовательного вычисления выходов каждого слоя сети от входов к выходу
\item Копирование сети на другие компьютеры
\item Удаление шума из данных
\end{itemize}
{\small\textit{Правильный ответ: Процесс последовательного вычисления выходов каждого слоя сети от входов к выходу.}}

\item Какое важное свойство должна иметь функция активации нейронов скрытого слоя, чтобы можно было применить алгоритм обратного распространения ошибки?
\begin{itemize}[label=$\square$, leftmargin=1.5cm, itemsep=0.2em, topsep=0.3em]
\item Быть ступенчатой (пороговой)
\item Быть дифференцируемой (иметь вычислимую производную)
\item Принимать только целочисленные значения
\item Быть строго убывающей
\end{itemize}
{\small\textit{Правильный ответ: Быть дифференцируемой (иметь вычислимую производную).}}

\item Какая проблема часто возникает при использовании сигмоиды в глубоких нейронных сетях?
\begin{itemize}[label=$\square$, leftmargin=1.5cm, itemsep=0.2em, topsep=0.3em]
\item Проблема затухающего градиента (Vanishing gradient)
\item Ошибка компиляции
\item Деление на ноль
\item Преобразование в целые числа
\end{itemize}
{\small\textit{Правильный ответ: Проблема затухающего градиента (Vanishing gradient).}}

\item В контексте аппроксимации функции, сколько нейронов должно быть в выходном слое для предсказания одного значения (одномерной функции)?
\begin{itemize}[label=$\square$, leftmargin=1.5cm, itemsep=0.2em, topsep=0.3em]
\item 0
\item 1
\item 2
\item Зависит от количества обучающих примеров
\end{itemize}
{\small\textit{Правильный ответ: 1.}}

\item Какую функцию активации чаще всего используют на ВЫХОДНОМ слое для задачи аппроксимации (регрессии)?
\begin{itemize}[label=$\square$, leftmargin=1.5cm, itemsep=0.2em, topsep=0.3em]
\item Сигмоиду
\item Softmax
\item Линейную (отсутствие активации)
\item Пороговую функцию Хевисайда
\end{itemize}
{\small\textit{Правильный ответ: Линейную (отсутствие активации).}}

\item Как в матричном виде вычисляется выход скрытого слоя H? (X — входы, W — веса, b — смещение, f — функция активации)
\begin{itemize}[label=$\square$, leftmargin=1.5cm, itemsep=0.2em, topsep=0.3em]
\item H \\textbackslash{}= X · W + b (без функции активации)
\item H \\textbackslash{}= f(X · W + b) (матричное произведение входов на веса, плюс смещение, затем функция активации)
\item H \\textbackslash{}= X · W (без смещения и без активации)
\item H \\textbackslash{}= f(W) + X
\end{itemize}
{\small\textit{Правильный ответ: H \\textbackslash{}= f(X · W + b) (матричное произведение входов на веса, плюс смещение, затем функция активации).}}

\item Что такое локальный градиент (дельта) выходного нейрона при линейной функции активации и функции потерь MSE?
\begin{itemize}[label=$\square$, leftmargin=1.5cm, itemsep=0.2em, topsep=0.3em]
\item Разность между предсказанным и истинным значением (ошибка)
\item Произведение предсказанного и истинного значений
\item Сумма всех весов сети
\item Всегда 0
\end{itemize}
{\small\textit{Правильный ответ: Разность между предсказанным и истинным значением (ошибка).}}

\item Как вычисляется локальный градиент (дельта) для нейрона скрытого слоя?
\begin{itemize}[label=$\square$, leftmargin=1.5cm, itemsep=0.2em, topsep=0.3em]
\item Он случайным образом генерируется на каждой эпохе
\item Как ошибка умноженная на производную функции активации скрытого слоя
\item Путем суммирования взвешенных дельт следующего слоя, умноженных на производную функции активации текущего слоя
\item Он равен выходу этого нейрона
\end{itemize}
{\small\textit{Правильный ответ: Путем суммирования взвешенных дельт следующего слоя, умноженных на производную функции активации текущего слоя.}}

\item Что произойдет, если все веса многослойной сети инициализировать одинаковыми значениями (например, нулями)?
\begin{itemize}[label=$\square$, leftmargin=1.5cm, itemsep=0.2em, topsep=0.3em]
\item Сеть обучится быстрее
\item Возникнет проблема симметрии: все нейроны будут обновляться одинаково и сеть не сможет выучить сложные закономерности
\item Сеть идеально аппроксимирует линейную функцию
\item Ошибка мгновенно станет нулевой
\end{itemize}
{\small\textit{Правильный ответ: Возникнет проблема симметрии: все нейроны будут обновляться одинаково и сеть не сможет выучить сложные закономерности.}}

\item Как обычно инициализируют веса в многослойном персептроне?
\begin{itemize}[label=$\square$, leftmargin=1.5cm, itemsep=0.2em, topsep=0.3em]
\item Нулями
\item Единицами
\item Малыми случайными числами
\item Бесконечно большими числами
\end{itemize}
{\small\textit{Правильный ответ: Малыми случайными числами.}}

\item Что означает термин "эпоха" при обучении нейросети?
\begin{itemize}[label=$\square$, leftmargin=1.5cm, itemsep=0.2em, topsep=0.3em]
\item Одно обновление весов
\item Один проход алгоритма прямого и обратного распространения по всей обучающей выборке
\item Одно предсказание
\item Время отклика сети
\end{itemize}
{\small\textit{Правильный ответ: Один проход алгоритма прямого и обратного распространения по всей обучающей выборке.}}

\item Какой параметр регулирует, насколько сильно изменяются веса на каждом шаге алгоритма Backpropagation?
\begin{itemize}[label=$\square$, leftmargin=1.5cm, itemsep=0.2em, topsep=0.3em]
\item Смещение (Bias)
\item Коэффициент скорости обучения (Learning rate / alpha)
\item Количество нейронов в скрытом слое
\item Производная
\end{itemize}
{\small\textit{Правильный ответ: Коэффициент скорости обучения (Learning rate / alpha).}}

\item Что такое переобучение (overfitting) в задаче аппроксимации функций?
\begin{itemize}[label=$\square$, leftmargin=1.5cm, itemsep=0.2em, topsep=0.3em]
\item Когда сеть плохо работает и на обучающей, и на тестовой выборках
\item Когда сеть идеально запоминает обучающую выборку со всеми шумами, но выдает плохие результаты на новых (тестовых) данных
\item Обучение сети слишком быстро
\item Когда функция активации переполняется
\end{itemize}
{\small\textit{Правильный ответ: Когда сеть идеально запоминает обучающую выборку со всеми шумами, но выдает плохие результаты на новых (тестовых) данных.}}

\item Как можно бороться с переобучением нейронной сети?
\begin{itemize}[label=$\square$, leftmargin=1.5cm, itemsep=0.2em, topsep=0.3em]
\item Увеличить количество нейронов до максимума
\item Убрать функцию активации
\item Использовать регуляризацию, раннюю остановку (early stopping) или dropout
\item Сделать скорость обучения равной 1
\end{itemize}
{\small\textit{Правильный ответ: Использовать регуляризацию, раннюю остановку (early stopping) или dropout.}}

\item Какое преимущество функции активации ReLU по сравнению с сигмоидой в скрытых слоях глубоких нейронных сетей?
\begin{itemize}[label=$\square$, leftmargin=1.5cm, itemsep=0.2em, topsep=0.3em]
\item Она ограничена значениями от 0 до 1
\item Она смягчает проблему затухающего градиента и ускоряет обучение за счёт вычислительной простоты
\item Она всегда дифференцируема в каждой точке
\item Она гарантирует отсутствие переобучения
\end{itemize}
{\small\textit{Правильный ответ: Она смягчает проблему затухающего градиента и ускоряет обучение за счёт вычислительной простоты.}}

\item Если матрица входных данных содержит 100 примеров по 5 признаков (размерность 100×5), а в скрытом слое 10 нейронов (матрица весов 5×10), какова размерность выхода скрытого слоя до применения функции активации?
\begin{itemize}[label=$\square$, leftmargin=1.5cm, itemsep=0.2em, topsep=0.3em]
\item 100×5 (число примеров на число признаков)
\item 5×10 (число признаков на число нейронов)
\item 100×10 (число примеров на число нейронов скрытого слоя)
\item 10×100 (число нейронов на число примеров)
\end{itemize}
{\small\textit{Правильный ответ: 100×10 (число примеров на число нейронов скрытого слоя).}}

\item Для задачи аппроксимации функции y \\textbackslash{}= sin(x) целесообразно ли нормировать входные данные?
\begin{itemize}[label=$\square$, leftmargin=1.5cm, itemsep=0.2em, topsep=0.3em]
\item Нет, нейросети работают с любым диапазоном без потери производительности
\item Да, масштабирование помогает избежать попадания в насыщение функций активации и ускоряет обучение
\item Нормирование полезно только для деревьев решений
\item Это приведет к удалению полезного сигнала
\end{itemize}
{\small\textit{Правильный ответ: Да, масштабирование помогает избежать попадания в насыщение функций активации и ускоряет обучение.}}

\item Какие данные используются при тестировании способности сети к обобщению аппроксимируемой функции?
\begin{itemize}[label=$\square$, leftmargin=1.5cm, itemsep=0.2em, topsep=0.3em]
\item Те же данные, что и для обучения
\item Генерация случайного шума
\item Отложенная тестовая выборка, которую сеть не видела при обучении
\item Только коэффициенты весов
\end{itemize}
{\small\textit{Правильный ответ: Отложенная тестовая выборка, которую сеть не видела при обучении.}}

\item Чем пакетный (batch) градиентный спуск отличается от стохастического (stochastic)?
\begin{itemize}[label=$\square$, leftmargin=1.5cm, itemsep=0.2em, topsep=0.3em]
\item Стохастический использует случайные функции активации
\item В пакетном обновляются веса один раз после просмотра всех примеров, а в стохастическом – после просмотра каждого отдельного примера
\item Пакетный использует только один слой сети
\item Отличий нет
\end{itemize}
{\small\textit{Правильный ответ: В пакетном обновляются веса один раз после просмотра всех примеров, а в стохастическом – после просмотра каждого отдельного примера.}}

\item Что такое функция потерь (Loss function)?
\begin{itemize}[label=$\square$, leftmargin=1.5cm, itemsep=0.2em, topsep=0.3em]
\item Функция, описывающая время работы скрипта
\item Функция, вычисляющая меру рассогласования между предсказанием сети и истинными ответами
\item Функция активации скрытого слоя
\item Функция распределения начальных весов
\end{itemize}
{\small\textit{Правильный ответ: Функция, вычисляющая меру рассогласования между предсказанием сети и истинными ответами.}}

\item Зачем нужно вычислять производную функции потерь по весам сети?
\begin{itemize}[label=$\square$, leftmargin=1.5cm, itemsep=0.2em, topsep=0.3em]
\item Для нахождения направления наискорейшего убывания ошибки
\item Для распечатки логов на экран
\item Для нормализации данных
\item Для увеличения числа нейронов
\end{itemize}
{\small\textit{Правильный ответ: Для нахождения направления наискорейшего убывания ошибки.}}

\item Какую форму принимает график функции потерь на обучающей выборке у успешно обучающейся нейронной сети?
\begin{itemize}[label=$\square$, leftmargin=1.5cm, itemsep=0.2em, topsep=0.3em]
\item Экспоненциально растет
\item Постоянно колеблется около нуля
\item Радикально возрастает в начале и падает в конце
\item Постепенно убывает, выходя на асимптоту (плато)
\end{itemize}
{\small\textit{Правильный ответ: Постепенно убывает, выходя на асимптоту (плато).}}

\item Достаточно ли одного нейрона в скрытом слое для качественной аппроксимации сложной нелинейной функции, имеющей несколько экстремумов?
\begin{itemize}[label=$\square$, leftmargin=1.5cm, itemsep=0.2em, topsep=0.3em]
\item Да, один нейрон может все
\item Нет, чем сложнее функция, тем больше может потребоваться нейронов или слоев для гибкости модели
\item Да, если использовать линейную активацию
\item Нет, для таких функций подходят только сети Кохонена
\end{itemize}
{\small\textit{Правильный ответ: Нет, чем сложнее функция, тем больше может потребоваться нейронов или слоев для гибкости модели.}}

\item Может ли алгоритм Backpropagation застрять в локальном минимуме функции потерь?
\begin{itemize}[label=$\square$, leftmargin=1.5cm, itemsep=0.2em, topsep=0.3em]
\item Строго нет, он всегда находит глобальный оптимум
\item Да, так как поверхность функции потерь в глубоких сетях невыпуклая
\item Да, но только при определённой скорости обучения
\item Нет, если скорость обучения нулевая
\end{itemize}
{\small\textit{Правильный ответ: Да, так как поверхность функции потерь в глубоких сетях невыпуклая.}}

\item Какая операция позволяет сдвигать каждый выходной нейрон скрытого слоя независимо?
\begin{itemize}[label=$\square$, leftmargin=1.5cm, itemsep=0.2em, topsep=0.3em]
\item Прибавление вектора смещений (bias)
\item Умножение на ноль
\item Транспонирование
\item Инверсия
\end{itemize}
{\small\textit{Правильный ответ: Прибавление вектора смещений (bias).}}

\end{enumerate}
